\documentclass{article}
\usepackage{microtype}
\usepackage{graphicx}
\usepackage{booktabs}
\usepackage{hyperref}
\newcommand{\theHalgorithm}{\arabic{algorithm}}
\usepackage{icml2026}
\usepackage{amsmath}
\usepackage{amssymb}

\icmltitlerunning{Predicting Cover Times on Small Graphs}

\begin{document}

\twocolumn[
\icmltitle{Simple Statistics Predict Random-Walk Cover Times on Small Graphs \\ Better Than the Universal Bound}

\icmlsetsymbol{equal}{*}
\begin{icmlauthorlist}
\icmlauthor{Anonymous}{anon}
\end{icmlauthorlist}
\icmlaffiliation{anon}{Autonomous research run; affiliation withheld}
\icmlcorrespondingauthor{Anonymous}{anon@example.org}
\icmlkeywords{random walks, cover time, Monte Carlo, small graphs}
\vskip 0.3in
]
\printAffiliationsAndNotice{}

\begin{abstract}
The cover time of a simple random walk is the expected number of steps needed to visit every vertex. For connected graphs with at most 30 vertices, the general bound $2m(n-1)$ is often loose. We ask whether a few cheap graph statistics (vertices $n$, edges $m$, diameter $D$, maximum degree $\Delta$) can predict cover time more accurately, and where such a predictor fails. We simulate $2000$ walks from random start vertices on $126$ connected graphs from $13$ families (paths, cycles, stars, complete graphs, wheels, grids, lollipops, barbells, trees, sparse and dense random graphs, random cubic graphs) with $n\in[8,28]$. The median ratio of cover time to $2m(n-1)$ is $0.195$. A log-linear model in $n,m,D,\Delta$, scored by leave-one-family-out cross-validation, has median absolute percentage error $0.14$, versus $0.64$ for the calibrated bound. The error is concentrated on dense graphs and grids. A replication with a second seed gives the same ordering.
\end{abstract}

\section{Introduction}
Cover times underlie randomized graph exploration, universal traversal sequences, and maze search~\cite{aleliunas1979}. Feige's sharp asymptotic analysis~\cite{feige1995} and surveys such as Lov\'asz~\cite{lovasz1993} characterize worst cases over all graphs. Worst-case constants, however, say little about a specific small graph. For $n\le 30$ the gap between the $2m(n-1)$ bound and the truth can exceed an order of magnitude on dense graphs (Table~\ref{tab:fam}). We ask a practical question: can a handful of easily computed statistics predict the cover time of a small connected graph more accurately than the general bound, and which graph types break the prediction? Our contribution is an empirical answer with uncertainty quantified by Monte Carlo and by out-of-family validation. We make no theoretical claim.

\section{Method}
\textbf{Graphs.} We generate 13 families (Table~\ref{tab:fam}) with target sizes $n\in\{8,12,16,20,24,28\}$. Deterministic families use one graph per size; random families (random trees, $G(n,0.12)$, $G(n,0.4)$ conditioned on connectivity, random cubic graphs via the configuration model with rejection) use three seeds per size. This yields $126$ connected graphs. Grids use $\lfloor\sqrt n\rfloor$ columns, so $n$ is rounded down.

\textbf{Cover time.} For each graph we run $W=2000$ simple random walks, each from a uniformly random start vertex, and record the first time all vertices have been visited. We report the sample mean $\hat C$ and a $95\%$ normal confidence half-width $1.96\,s/\sqrt W$. The largest relative half-width over all graphs is $5.4\%$. As a sanity check, simulated values match known closed forms: the cycle gives $376.9$ at $n=28$ against $n(n-1)/2=378$, and $K_{28}$ gives $104.3$ against $(n-1)H_{n-1}\approx 104.4$.

\textbf{Predictors.} Let $L=\log$. We compare four models, each fitted by ordinary least squares on $\log \hat C$:
\begin{itemize}
\item \texttt{bound}: $\log \hat C = \log(2m(n-1)) + \kappa$, with one fitted constant $\kappa$. This is the calibrated general bound.
\item \texttt{nm}: $\log \hat C = \beta_0 + \beta_1\log n + \beta_2\log m$.
\item \texttt{nmD}: adds $\beta_3 \log D$.
\item \texttt{nmDd}: adds $\beta_4\log\Delta$.
\end{itemize}
Each model is evaluated by leave-one-family-out (LOFO) cross-validation. For each family we train on the other 12 and predict the held-out family. We report the median and mean absolute percentage error (APE), $|\hat C_{\text{pred}}-\hat C|/\hat C$, pooled over all graphs.

\section{Experiments}
\begin{table}[t]
\caption{Median leave-one-family-out APE per family (seed 0, $W=2000$). Columns are predictors. Lower is better. Rows are ordered by the \texttt{bound} column, from least to most affected.}
\label{tab:fam}
\begin{center}
\small
\begin{tabular}{lrrrr}
\toprule
Family & \texttt{bound} & \texttt{nm} & \texttt{nmD} & \texttt{nmDd} \\
\midrule
path & 0.751 & 0.477 & 0.114 & 0.015 \\
cycle & 0.404 & 0.170 & 0.466 & 0.047 \\
barbell & 0.471 & 0.586 & 0.244 & 0.042 \\
regular3 & 0.739 & 0.819 & 0.711 & 0.074 \\
gnp\_sparse & 0.471 & 0.120 & 0.096 & 0.072 \\
wheel & 1.002 & 0.473 & 0.265 & 0.101 \\
random\_tree & 0.664 & 0.188 & 0.222 & 0.129 \\
lollipop & 0.611 & 0.817 & 0.662 & 0.458 \\
binary\_tree & 0.625 & 0.196 & 0.217 & 0.343 \\
star & 0.242 & 0.676 & 0.590 & 0.191 \\
complete & 14.601 & 0.209 & 0.347 & 0.569 \\
grid & 0.382 & 0.417 & 0.729 & 0.568 \\
gnp\_dense & 3.674 & 0.447 & 0.424 & 0.578 \\
\midrule
\textbf{All (median)} & \textbf{0.635} & \textbf{0.350} & \textbf{0.301} & \textbf{0.139} \\
\textbf{All (mean)} & 1.922 & 0.428 & 0.363 & 0.236 \\
\bottomrule
\end{tabular}
\end{center}
\vskip -0.1in
\end{table}

\textbf{Results.} Across the 126 graphs the simulated cover time is a median $0.195$ of $2m(n-1)$, ranging from $0.005$ to $0.589$. The universal bound is therefore loose by a factor of about five at the median and by about two hundred on the densest graphs. The calibrated bound model fails badly on the complete and dense random graphs (Table~\ref{tab:fam}), because its single constant is set by sparse graphs. The $\log n,\log m$ model already halves the median error. Adding the diameter lowers it further, and adding $\Delta$ brings the median LOFO APE to $0.139$. The second seed gives $0.626$, $0.322$, $0.286$, and $0.144$ for the four models, so the ordering is stable. The full-data fit of \texttt{nmDd} is
\[
\hat C \approx e^{-0.61}\, n^{1.06}\, m^{-0.10}\, D^{1.14}\, \Delta^{0.59}.
\]
The exponents of $n$ and $D$ are both near one, and $m$ contributes little once $D$ and $\Delta$ are known. This suggests that the dominant structure is path-like, which is consistent with the cover time being controlled by slow spreading along long paths and bottlenecks. We do not interpret the coefficients causally.

\section{Discussion and Limitations}
\textbf{Where the predictor is reliable.} The \texttt{nmDd} model is accurate, with median APE under $0.1$, on path-like and sparse graphs: paths, cycles, barbells, sparse $G(n,p)$, and cubic graphs. There the cover time is governed by $n$, $D$, and bottleneck structure.

\textbf{Where it fails.} The error is largest on dense and highly symmetric graphs: complete graphs ($0.57$), dense $G(n,0.4)$ ($0.58$), and grids ($0.57$). For $K_n$ the true cover time is $\Theta(n\log n)$ with $D=1$, which lies outside the range of the other families, so LOFO extrapolates poorly. Lollipops are also hard ($0.46$) because their cover time depends on where the walk starts relative to the clique, a feature we do not measure.

\textbf{Limitations.} (i) The sample is small: 13 families, sizes at six points, and one graph per size for deterministic families. LOFO error on a family depends on how much its neighbours resemble it. (ii) The start vertex is uniform, and we average over it, so the predictor targets the average, not the worst case. (iii) The Monte Carlo error is small ($\le 5.4\%$ half-width), but it is far below the model error, so it does not limit the conclusions. (iv) We compare against only one general bound. We did not implement sharper worst-case constants. (v) The results are for $n\le 28$ and do not necessarily extend to large graphs.

\section{Reproducibility}
Code is in \texttt{code/}: \texttt{exp.py} generates graphs, runs the walks, and writes JSON; \texttt{analyze.py} performs the fitting and LOFO evaluation. Pure Python 3, standard library only. Commands: \texttt{python3 -I exp.py --seed 0 --walks 2000 --out results.json}, then \texttt{python3 -I analyze.py results.json}. The replicate uses \texttt{--seed 1}. A run takes about 1 minute on one CPU core. Walk randomness and graph generation both come from \texttt{random.Random(seed)}.

\begin{thebibliography}{9}
\bibitem[Aleliunas et~al.(1979)]{aleliunas1979}
R.~Aleliunas, R.~M. Karp, R.~J. Lipton, L.~Lov\'asz, and C.~Rackoff.
\newblock Random walks, universal traversal sequences, and the complexity of maze problems.
\newblock In \emph{FOCS}, 1979.
\bibitem[Feige(1995)]{feige1995}
U.~Feige.
\newblock A tight upper bound on the cover time for random walks on graphs.
\newblock \emph{Random Structures \& Algorithms}, 1995.
\bibitem[Lov\'asz(1993)]{lovasz1993}
L.~Lov\'asz.
\newblock Random walks on graphs: a survey.
\newblock In \emph{Combinatorics, Paul Erd\H{o}s is Eighty}, 1993.
\end{thebibliography}

\end{document}
